HeatCast-FM shows that pretrained time-series foundation models can serve as the forecasting core of a regional heatwave warning system without any training on Indian data. Used zero-shot on IMD gridded Tmax, Chronos-2 with covariates gave 7-day probabilistic forecasts about 17\% better than climatology on CRPS. Its hot-day probabilities were reliable, and its heatwave ROC area was above 0.96. Models trained on 64 years of local data were still about 6\% better, with the gap growing with lead time. On this evidence, a foundation model is a good starting point for places or variables with short records. Where long local records exist, a well-built statistical or LSTM model remains hard to beat.

For the use case, the work delivers:
\begin{itemize}
\item a reproducible pipeline from IMD data to colour warnings;
\item a way to turn any quantile forecast into IMD rule probabilities for both coastal and plains stations;
\item a tested point correction that roughly halved the point error at Mumbai and Ratnagiri;
\item evidence that the Tmax-only rule misses many humid heat days on the Konkan coast;
\item a small viewer app that shows the 7-day forecast, event probabilities and warning level for any region and date.
\end{itemize}

Future work:
\begin{itemize}
\item replace the station proxy with the planned IoT weather stations and repeat the point validation;
\item complete the LoRA fine-tuning of Chronos-2 on a GPU;
\item calibrate the Chronos intervals on validation years, since they are too narrow;
\item test station data or finer temperature grids, where the official heatwave rule can be applied as intended;
\item add a humidity-aware warning level once IMD publishes its revised criteria.
\end{itemize}
